\section{Introduction}
\label{sec:introduction}

Neural scaling laws have ushered in new paradigms for interpreting deep learning dynamics. These laws relate performance gains to model size, dataset scale, and compute, motivating longer training horizons under tight efficiency constraints~\cite{hestness2017deep,kaplan2020scaling,henighan2020scaling,hoffmann2022chinchilla}. In this regime, optimization stability becomes an important practical requirement, since unstable training may result in slow convergence and high sensitivity to hyperparameter choices. In practice, stability is often maintained through a set of heuristics, such as momentum, batch-size scaling, normalization, weight decay, and learning-rate schedules~\cite{goyal2017accurate,smith2017cyclical,loshchilov2016sgdr,ioffe2015batch,krogh1991simple}. Despite their effectiveness, understanding when these mechanisms help, and when they fail, remains consequential for the design of training procedures.

Training can become unstable when updates are too large relative to sharp directions of the loss landscape, so that they overshoot rather than converge toward a minimum~\cite{polyak1964methods,bertsekas1997nonlinear}. Stochastic extensions, such as the break-even point framework, relate instability to both sharpness and curvature variability, particularly during an early phase in which learning rate and batch size shape the subsequent trajectory~\cite{jastrzebski2020break}. The Edge-of-Stability literature, however, shows that deep networks often continue to make progress near, or even beyond, the nominal sharpness threshold~\cite{cohen2021edge,arora2022understanding,damian2023self}. These works provide an important geometric interpretation of training dynamics. They also motivate a complementary question: how does one update interact with the next, and how do their repeated interactions affect learning? A trajectory can repeatedly reverse in some directions while continuing to descend in others. Explaining this behavior requires following both the magnitude and direction of the steps actually taken, rather than treating every oscillation as failed progress.

To examine these interactions, we characterize constructive and destructive training regimes through consecutive-gradient alignment and gradient magnitude. Persistent positive alignment, $\rho_t>0$, identifies a constructive regime in which successive SGD steps reinforce one another, whereas persistent negative alignment, $\rho_t<0$, identifies a destructive regime in which successive steps partially reverse preceding parameter changes. We connect these regimes to directional curvature and introduce an inexpensive estimate of curvature along the step taken. It uses consecutive gradients and the actual parameter change, so the formulation applies to SGD, Adam, and other optimizers without additional Hessian computations. Under plain SGD, it reduces to a combination of gradient norms and their inner product. The formulation also makes the contribution of mini-batch sampling explicit, clarifying how noise affects the observed alignment and curvature estimate.

Destructive alignment alone does not imply instability. Within a local quadratic approximation, reversal and gradient growth jointly determine whether an SGD step decreases or increases the loss. This relationship explains how destructive regimes can sustain oscillatory descent, while sufficiently amplified reversals can produce upward loss excursions. Reading alignment and gradient norms together therefore provides a geometric interpretation of signals already available during training. Normalized displacement complements this local view by measuring how much of the total distance traveled remains as net parameter change.

We examine these relationships through controlled full-batch and mini-batch experiments across image-classification and language-modeling tasks, using multilayer perceptrons, residual networks, and Transformers. The SGD experiments isolate curvature and sampling effects. The Adam experiments then follow the argument from consecutive-gradient behavior to loss spikes, from spikes to changes in the network's representations, and from those changes to interventions. Evaluating the same reference examples at successive training steps makes it possible to see where representations reverse and where their changes become large. This motivates a local displacement cap, which we compare with learning-rate reduction, second-moment changes, and global gradient clipping from the same training states.

Our central contribution is to connect the interaction between training
steps to its consequences in parameter space, in the loss, and in the
learned function. This connection provides both an interpretation of
oscillatory progress and a measurable basis for intervention. Specifically:
\begin{itemize}
\item \textbf{Relating local reversal to accumulated learning dynamics.}
Edge-of-Stability studies establish that oscillation can coexist with
progress~\cite{cohen2021edge,damian2023self}, while gradient-product
diagnostics identify transitions toward
stationarity~\cite{chee2018diagnostics,pesme2020diagnostics}. These findings
motivate distinguishing opposition between neighboring updates from
cancellation across the complete trajectory. We connect consecutive-gradient
alignment to an exact all-pair displacement decomposition and normalized
displacement. The GD experiments show that negative adjacent interactions
can coexist with positive total interaction and decreasing loss, explaining
why a destructive local regime need not imply failed progress
(Sections~\ref{sec:displacement}--\ref{sec:gd-experiments}).

\item \textbf{Measuring curvature along the step actually taken.}
The leading Hessian eigenvalue used in spectral stability
analysis~\cite{cohen2021edge} describes the sharpest direction, but does not
specify the curvature encountered by a particular optimizer update.
We introduce an inexpensive step-direction estimate from consecutive
gradients and the applied parameter change, without additional backward
passes or Hessian computations. For smooth full objectives, we establish
its exact step-average interpretation independently of the optimizer,
including Adam, and bound its starting-point approximation error. Its
SGD specialization exposes the roles of alignment, norm growth, and
sampling noise; independently computed Hessian--vector products test
that interpretation (Sections~\ref{subsec:curvature-diagnostic}
and~\ref{subsec:gd-validation}).

\item \textbf{Separating gradient reversal from loss-spike formation.}
Bai et al.~\cite{bai2026spikes} explain Adam spikes through adaptive
preconditioning and gradient-directional curvature. To connect spike
behavior to consecutive training gradients, we express the local SGD
loss-increase condition jointly in terms of reversal and norm growth,
then examine those quantities under Adam without transferring the SGD
threshold into an Adam guarantee. An exhaustive comparison across
loss-increase magnitudes and future steps identifies both the association
and its missed increases and false alarms, extending the characterization
beyond selected peaks (Sections~\ref{subsec:loss}
and~\ref{subsec:characterization}).

\item \textbf{Tracing functional oscillation and testing selective control.}
Representation-change correction~\cite{kosson2024warmup} and unit-wise
gradient clipping~\cite{brock2021normalization} already motivate controlling
update size. We examine where successive representation changes reverse,
which patterns carry their movement, and whether limiting that movement
improves an oscillatory trajectory. Tracing the same reference examples
reveals concentrated movement and shared temporal structure across
measured stages. We then cap the direct affine displacement induced by
Adam's proposal using an absolute per-unit budget. Matched comparisons
with learning-rate reduction, second-moment changes, and global clipping
show lower mean training loss and positive loss variation for all tested
cap quantiles in the three cases, while reversals persist. Endpoint,
held-out, and cost tradeoffs delimit this training-dynamics result
(Sections~\ref{subsec:tracing}--\ref{sec:vc-interventions}).
\end{itemize}

The remainder of the paper is organized as follows. Section~\ref{sec:background} reviews training-stability perspectives. Section~\ref{sec:displacement} develops the displacement formulation and its diagnostics. Section~\ref{sec:gd-experiments} examines full-batch GD and mini-batch sampling. Section~\ref{sec:adam} develops the Adam characterization, representation analysis, and interventions. Section~\ref{sec:conclusion} concludes.
